% Options for packages loaded elsewhere
\PassOptionsToPackage{unicode}{hyperref}
\PassOptionsToPackage{hyphens}{url}
\documentclass[
]{article}
\usepackage{xcolor}
\usepackage{amsmath,amssymb}
\setcounter{secnumdepth}{-\maxdimen} % remove section numbering
\usepackage{iftex}
\ifPDFTeX
  \usepackage[T1]{fontenc}
  \usepackage[utf8]{inputenc}
  \usepackage{textcomp} % provide euro and other symbols
\else % if luatex or xetex
  \usepackage{unicode-math} % this also loads fontspec
  \defaultfontfeatures{Scale=MatchLowercase}
  \defaultfontfeatures[\rmfamily]{Ligatures=TeX,Scale=1}
\fi
\usepackage{lmodern}
\ifPDFTeX\else
  % xetex/luatex font selection
\fi
% Use upquote if available, for straight quotes in verbatim environments
\IfFileExists{upquote.sty}{\usepackage{upquote}}{}
\IfFileExists{microtype.sty}{% use microtype if available
  \usepackage[]{microtype}
  \UseMicrotypeSet[protrusion]{basicmath} % disable protrusion for tt fonts
}{}
\makeatletter
\@ifundefined{KOMAClassName}{% if non-KOMA class
  \IfFileExists{parskip.sty}{%
    \usepackage{parskip}
  }{% else
    \setlength{\parindent}{0pt}
    \setlength{\parskip}{6pt plus 2pt minus 1pt}}
}{% if KOMA class
  \KOMAoptions{parskip=half}}
\makeatother
\ifLuaTeX
  \usepackage{luacolor}
  \usepackage[soul]{lua-ul}
\else
  \usepackage{soul}
\fi
\setlength{\emergencystretch}{3em} % prevent overfull lines
\providecommand{\tightlist}{%
  \setlength{\itemsep}{0pt}\setlength{\parskip}{0pt}}
\usepackage{bookmark}
\IfFileExists{xurl.sty}{\usepackage{xurl}}{} % add URL line breaks if available
\urlstyle{same}
\hypersetup{
  pdftitle={Examplar 1},
  hidelinks,
  pdfcreator={LaTeX via pandoc}}

\title{Examplar 1}
\author{}
\date{}

\begin{document}
\maketitle

\section{Abstract 1:}\label{abstract-1}

Models derived from physical laws (deductive models) and those obtained
from experimental observations and operational data (inductive) are
widely used in mechanical engineering to support operational
optimisation and structural integrity assessment. However, Digital Twins
are still commonly described as either physics-based or data-driven
systems, masking the hybrid nature that progressively emerges during
their development, deployment, and evolution. Although hybrid modeling
is frequently implemented in practice, it remains insufficiently
characterized, preventing engineers from intentionally designing hybrid
Digital Twins.

This paper analyzes two mechanical engineering applications---a fluidic
loop optimized through predictive operation and a pressure vessel
Digital Twin dedicated to lifetime estimation---to investigate how
hybridization appears during model creation, operational usage, and
model updating stages. The study identifies when and why hybridization
occurs across Digital Twin engineering activities.

Results show that hybridization systematically emerges regardless of the
initial modeling paradigm(building inductives or deductive models).
\st{The interaction between deductive and inductive models enables
adaptive prediction capabilities and continuity across lifecycle
phases.}

Recognizing Digital Twins as intrinsically hybrid systems provides a
framework for more resilient and lifecycle-oriented engineering
processes. Hybrid modeling should therefore be considered a foundational
principle rather than an optional enhancement, enabling transferable
methodologies for future Digital Twin developments.

\subsubsection{1. Introduction}\label{introduction}

\begin{itemize}
\tightlist
\item
  DT definition in mechanical engineering
\item
  Physics vs data-driven dichotomy
\item
  Research gap: lack of explicit hybrid characterization
\item
  Contributions of the paper
\end{itemize}

\subsubsection{2. Hybridization in DT within Engineering Approach and
Case
Studies}\label{hybridization-in-dt-within-engineering-approach-and-case-studies}

\begin{itemize}
\tightlist
\item
  Definitions:

  \begin{itemize}
  \tightlist
  \item
    Model
  \item
    Model nature
  \item
    Hybridization process
  \item
    Model lifecycle
  \end{itemize}
\item
  Fluidic loop DT

  \begin{itemize}
  \tightlist
  \item
    objective
  \item
    DT construction (models creation, verification, validation)

    \begin{itemize}
    \tightlist
    \item
      Hybridization process witin models lifecycle
    \end{itemize}
  \end{itemize}
\item
  Pressure vessel DT

  \begin{itemize}
  \tightlist
  \item
    lifetime estimation strategy
  \item
    DT constructions (models creation, verification, validation)

    \begin{itemize}
    \tightlist
    \item
      Hybridization process witin models lifecycle
    \end{itemize}
  \end{itemize}
\end{itemize}

\subsubsection{3. Comparative analysis between both
cases}\label{comparative-analysis-between-both-cases}

\subsubsection{4. Lessons Learned \&
Transferability}\label{lessons-learned-transferability}

\subsubsection{5. Conclusion}\label{conclusion}

\end{document}
